% --- Capítulo 8: Diseño de la red implementada ---
\chapter{Diseño de la Red Implementada}

\section{Arquitectura Global del Sistema de Fusión Triple}
% Descripción del "Jefe Final": cómo las tres ramas se integran en un único grafo computacional.

\section{Rama Visual Híbrida: Extractor RGB (CNN-ViT)}
\subsection{Módulo Convolucional: ResNet18 como extractor de rasgos locales}
\subsection{Módulo de Atención: Vision Transformer para relaciones globales}

\section{Rama de Análisis Estructural: Red Convolucional ARP}
% Diseño de la CNN específica para procesar proyecciones polares y bordes de la lesión.

\section{Rama de Percepción Clínica: Perceptrón Multicapa (MLP)}
% Implementación de la red para los Planes A y B de metadatos.

\section{Integración y Toma de Decisión: \textit{Late Fusion}}
% La capa de concatenación y el diseño del vector de características de 1280 dimensiones.
\subsection{Topología de las Cabezas de Decisión (Dual-Head)}

\section{Modelado Matemático del Entrenamiento}
\subsection{Optimización frente al desbalanceo: \textit{Focal Loss}}
\subsection{Función de pérdida combinada y ponderación de tareas}

\section{Tácticas de Entrenamiento Avanzado}
\subsection{Protocolo \textit{Linear Probing} y \textit{Fine-Tuning} global}
\subsection{Aprendizaje Jerárquico mediante \textit{Differential Learning Rate}}
% Explicación de las velocidades /100, /10 y /1 para cada sección de la red.
\subsection{Regularización por \textit{Modality Dropout}}
\subsection{Gestión de memoria y \textit{Gradient Accumulation}}
% Gradient Accumulation: Explica cómo lograste simular un Batch Size de 128 matemáticamente usando sub-lotes físicos de 32 (que sí cabían en tus 6GB de VRAM), acumulando gradientes en 4 pasos antes de actualizar los pesos.

% Automatic Mixed Precision (AMP): Detalla cómo activaste los Tensor Cores de las gráficas RTX/T4 para calcular en 16 bits (FP16) en lugar de 32 bits (FP32). Explica que esto no solo redujo a la mitad la memoria RAM de vídeo necesaria, sino que aceleró masivamente el cálculo matricial.

% Data Pipeline (Workers y Prefetch): Añade un párrafo explicando cómo, para evitar el GPU Starvation, configuraste el DataLoader explotando al máximo la CPU multihilo (num_workers=4, pin_memory=True y prefetch_factor=4), asegurando que la gráfica siempre tuviera datos listos para procesar.
\section{Optimización Bayesiana de Hiperparámetros (Optuna)}
% Espacio de búsqueda y configuración de las pruebas (trials).